% Faithful conversion of ages-i-sprint-4; original XML and source characters retained.
\hypertarget{__RefHeading___Toc4751_3118176633}{}
\refstepcounter{subsection}
\subsection*{2.3.5 Sprint 4}
\label{unit:ages-i-sprint-4}
\addcontentsline{toc}{subsection}{2.3.5 Sprint 4}

% text-p-0592
\indent \originaltab{}Inicialmente, para a última Sprint, reduzimos um pouco o escopo das tarefas a serem entregues devido ao curto tempo de desenvolvimento que tivemos na Sprint anterior. Porém, ao longo da semana, concluímos boa parte das tasks estipuladas. Assim sendo, adicionamos algumas funcionalidades a mais para serem entregues, como a opção de adicionar os suplementos escolhidos em um carrinho para serem enviados ao e-mail do usuário.

% text-p-0593
\indent \originaltab{}Com mais tempo de desenvolvimento em comparação com a Sprint 3, pude desenvolver algumas tasks mais trabalhosas. Realizei a tarefa de criar uma conexão entre todos os suplementos disponíveis no Banco de Dados com suas respectivas imagens. Para isso, tive de alterar o S\textit{chema} do Prisma para incluir um campo adicional chamado \textit{“path”}, que armazena o caminho para a imagem dos suplementos. Também tive de alterar os \textit{hooks} da loja e dos suplementos para lidarem com a nova lógica de vinculação de imagens. Ao final, tive de adicionar manualmente os \textit{links }das imagens para cada suplemento na \textit{seed }(dados de inicialização).

% text-p-0594
\indent \originaltab{}Também tive participação na tarefa de envio do carrinho (suplementos escolhidos) por e-mail, que realizei juntamente do meu colega Urien Nolasco (AGES I). Fui responsável por criar a estrutura capaz de receber a lista de suplementos a serem enviados ao usuário e também adicionei um \textit{loading }para aprimorar a \textit{User Experience}.

% text-p-0595
\indent \originaltab{}Para finalizar minha participação na AGES como AGES I, realizei alguns \textit{hotfixes} sugeridos pelo meu professor-orientador, Daniel Callegari, durante seu teste da aplicação. Para isso, tive apoio dos meus colegas (AGES III) Arthur Maia e (AGES I) Henrique Knack.

\begin{reflection}
% text-p-0606
\indent \originaltab{}Meus aprendizados durante a Sprint 4 foram muito significativos. Com a utilização da ORM Prisma, muitas tarefas e conceitos de Banco de Dados que eu acreditava serem complexas, se mostraram simples e práticas. Tive contato direto com o nosso Banco durante essa Sprint, e pude compreender ainda mais sobre adições e alterações na \textit{seed} e \textit{Data Base} em si. Essa tecnologia se mostrou tão útil e direta, que irei sugerir o uso da mesma em minha próximas participações da AGES.\originalspace{}

% text-p-0607
\indent \originaltab{}Em relação à equipe, todos trabalharam de maneira impecável para entregar um projeto robusto e completo aos nossos \textit{StakeHolders}, que, diga-se de passagem, adoraram o nosso trabalho.

\end{reflection}
